% 第1章 Introduction (原书页 1–8 / PDF p016–p023)
\chapter{引言}

\section{准备与教材}

这些讲义的主题是“固体理论中的概念”（Concepts in the Theory of Solids）——事实上我本该说“固体的\emph{量子}理论”，因为我们对物质性质的理解，几乎没有哪一点不在某种程度上依赖于量子理论。我们需要对量子理论有若干了解：即对矩阵形式体系（matrix formalism）与变换理论，以及对初等波动力学有相当程度的理解。例如，含时微扰论与算符运动方程将不再多加解释地直接使用；但现代场论的技术则不会这样，如有必要，我们会从头推导那些技术。

在固体领域需要何种准备，取决于本课程的意图：本课程不打算罗列固体性质的各种唯象表现，而是要稍微深入地探究这些性质背后的东西。在许多情形下，这意味着我们将试图理解固体\emph{为什么}如此行为；但在另外许多情形下——也许更多——我们将恰恰走到我们真正的理解失效之处。因此，对固体有哪些性质具备相当广泛的知识，显然会大有帮助。Kittel 的《Introduction to Solid State Physics》(1) 是一本出色的教科书，它在本课程所要讲述的内容的准备层次上纵览了整个领域。换言之，我们将假定读者对 Debye $T^3$ 比热、布里渊区（Brillouin zones）、自由电子比热或自旋顺磁性、电子或核顺磁共振，以及其他多少算得上标准的理论观念与实验技术已有一定的熟悉。自 Seitz 于 1940 年写成《Modern Theory of Solids》(2) 以来，还没有一本书试图在任一基础层次上完整覆盖固体理论；而且它事实上至今仍是远胜群伦的最好教材。这或许表明，自 1940 年以来我们在基础理解上并没有多大进展，只是在研究宽得多的现象类别——这在某种程度上是对的。无论如何，到目前为止，人们为把这部《现代理论》加以“现代化”所找到的唯一办法，就是出版一系列由综述文章组成的书，即所谓的“Seitz-schrift”或 Seitz-Turnbull 丛书（\emph{Solid State Physics---Advances in Research and Applications}），它大概是当前这一层次上最好的单一文献来源。Kittel 的书非常好，但除第二版中晶体对称性这一个领域可能是例外，它对任何一个专题都不算完备。某些专门领域在某些书里覆盖得相当不错——例如 Ziman 的《Electrons and Phonons》(3)，以及关于位错、磁共振等专门题目的其他书。还应提到 Wannier 的《Elements of Solid State Theory》(4) 与 Peierls 的书 (5)；这两本书各自大概都是某种选材的产物——类似于我们在这里将要做的选材，只是选法颇为不同。一部宏伟的、但相当高深又相当浓缩的教材，是 Landau 与 Lifshitz 的《Statistical Physics》(6)。即将问世的一本书将从群论与对称性原理的观点覆盖固体物理一个非常宽的领域，它就是 M.~Lax 的《Symmetry Principles in Solid-State Physics》(6a)。无论如何，在任一给定领域中用作资料来源的书与文章都将予以提及。

\section{课程计划}

\begin{enumerate}[label=\arabic*.,leftmargin=2.5em,itemsep=0.4em,parsep=0pt]
\item 先讲一些关于固体物理的一般观念，包括把固体分为各种类型的分类，以及关于我们或许希望从更深的观点加以解释的量子化学事实的一些相当宽泛的论述：不同类型固体的出现、性质等等。

\item 其次一个、也许是最重要也最基本的领域，是纯粹的单电子能带理论（one-electron band theory）；因为诸如结合、对称性、能带结构等等基本问题，可能多半都首先由能带决定。

  通过对较早的方法与结果作若干研究，我们希望最终引向的，是 Phillips、Cohen、Heine 等人关于近自由电子模型何以成功的一些较新的想法，以及在彼基础上人们可以就结合能的变化趋势与固体的量子化学所作的某些推测。

  另一个论题将是有微扰、外场、杂质等存在时单电子能带理论的近代进展，即有效质量理论（effective mass theory）。

\item 接着我们要评论：像上面讨论的单电子理论这样明显的过度简化为何竟能奏效——原因即在于元激发（elementary excitations）的观念，它大概是整个固体物理中单一最有成果的理论概念。在讨论了理论上简单的绝缘体情形与不那么简单的金属情形之后，我们将进而讨论集体元激发的各种可能形式——激子（exciton）、自旋波（spin wave）与声子（phonon）。还有一些适用于以上所有情形的关于集体激发的更一般的评述。

\item 在讨论自旋波时我们将处理磁性。这里的基本问题是：材料为何、何时具有磁性——即“磁状态”（magnetic state）的问题，或者说自由自旋何时出现——以及是什么引起自旋之间的相互作用并导致铁磁或反铁磁性。我们还将以此作为关于凝聚的若干一般性事实的一个例子。
\end{enumerate}

\section{固体的一般性质与分类}

我想，任何一门固体理论课程都应当从固体的定义讲起，尽管我们都知道，物理学家拒绝把固体定义为（粗略地说）踢上去弄痛你脚趾的那种东西，而是把它定义为原子的一种规则排列，其意义是：在某一空间群（space group）之一的操作下，它在很好的近似上具有平移对称性（translational symmetry）。

讲到这里，我们其实已经悄悄滑过了固体物理最根本的三个问题：（1）为什么存在固体？（2）如何从一种真正基本的观点——即原子核本身与电子都被真正按量子力学处理——来描述固体？（3）固体是靠什么、为什么把自己结合在一起的？

第一个问题可以成为激烈的语义学与哲学争论的话题，但据我所知，还没有人能给出一个真正的证明，甚至一个很好的定性理由，来说明为什么几乎所有原子集合体的基态（ground state）即最低能量态在本性上是规则的而非不规则的。例如，在最简单的可能情形——一个装满刚性球并受压的箱子——中，人人都出于显然的理由假定规则的密堆积（立方或六角）具有最小体积，因而在受压时能量最低；但据我所知，没有一个证明表明事情确实如此。

Landau 曾表述过这样一种观点 (7)：这类系统的基态必须具有某种不变性群。毕竟，初始的哈密顿量 $H_{\mathrm{el}}+H_{\mathrm{nuc}}$ 在完全的空间平移与转动群下是不变的。在我们知道的一个情形——低压下的 $\mathrm{He}_4$——以及可能还有另一个情形——$\mathrm{He}_3$——系统作为整体的基态保持了完全的平移与转动群，尽管在后一情形中得以维持的可能只是完全的平移对称性。在所有其他情形中，系统都会凝聚，我们的意思是它选择了一个还要更低的对称性，即一个周期性空间群。事实上，指望一个从具有完全平移—转动群对称性的哈密顿量出发的系统，其最低态竟全然没有任何对称性，是相当不合理的。也许更有趣的问题倒是：为什么量子液体不更多一些，而不是为什么没有类玻璃最低能量态的情形。我觉得这是一种有趣的看法，但很难算是一个令人信服的论证。

第二个问题，鉴于当前对固态 $\mathrm{He}_3$ 与 $\mathrm{He}_4$ 的兴趣，也许相当快就会至少以一种近似的方式得到解决；尽管如此，我还从未见过一个令人满意的、完全量子力学式的固体描述——从对原子间相互作用的真实描写出发，并把所有零点运动都适当地包括进来。

至于第三个问题——固体靠什么、为什么把自己结合在一起——我想，等我们讲到声子与集体激发时，会给出一个虽属近似却足够满意的说明。

于是，把固体确实存在当作一个实验事实，我们就可以一般地问：固体有哪些种类，以及可以如何对它们分类。对固体进行分类有若干种唯象的方法——例如，你们也许最熟悉的一种是按对称性分类，那是一种严格的数学处理方式，也非常有用，但与固体的结合力及其他物理性质没有直接关系。

在 Seitz 与 Kittel 的开篇论述中，可以找到一种与本课程要点切合得多的分类，它按照化学键的类型来分——从唯象的观点看，正是化学键负责着物质的结合能。Seitz 的分类包含五个类别：（1）金属；（2）离子晶体（ionic crystal）；（3）价晶体或共价晶体（valence or covalent crystal）；（4）分子晶体（molecular crystal）；（5）半导体。

如今我们已经认识到，从任何实际的意义上说，半导体与金属或价晶体之间在结合类型上的区别，以及半导体与任何其他类型绝缘体之间在导电性上的区别，都完全是人为的；半导体在任何实际意义上都不代表一类独立的晶体。Kittel 在其余四个类别之外增加了第五类：氢键型晶体（hydrogen-bonded crystal），从介电体与铁电体的角度看它意义重大，但在其他任何重要方面，它与分子晶体或离子晶体并没有什么不同。这样我们就得到了一个合理的五类分类（见表~1）。

\begin{table}[htbp]
\centering
{\bfseries 表~1}\\[6pt]
\small
\setlength{\tabcolsep}{3.5pt}
\renewcommand{\arraystretch}{1.2}
\begin{tabular}{@{}p{1.35cm}p{1.75cm}p{3.5cm}p{3.85cm}p{3.9cm}@{}}
\toprule
类型 & 类型或“典型”例子 & 结合的起因 & 性质 & 出现 \\
\midrule
分子晶体 & Ne 或 N$_2$ & 范德瓦耳斯吸引：相互极化，一种 $n$ 体效应 & 电阻率 $R$ 高。结合弱。$n^2$（折射率）与 $\epsilon$（介电常数）低：电子很少可动。密堆积结构 & (a) 稀有气体；(b) 周期表靠下、靠右的元素：C、N、O、H、F、S、P、Cl 等。\emph{无}电正性元素，重元素\emph{很少} \\
氢键型晶体 & H$_2$O & 氢键——O—H—O 使质子的动能降低 & 结合力中等偏弱（约为分子晶体的 10--100 倍）；$n^2$ 低，因质子可动而有奇特的介电性质。结构疏松 & 同上，但总含有 H 和一种电负性元素：O、F 等 \\
离子晶体 & NaCl & 离子键：Madelung 势 $+$ 使电离能取值恰当的壳层效应 & 结合力极强，结构颇为密堆积；$R$ 高，$n^2$ 低，$\epsilon$ 高（离子极化）；离子电导，低迁移率的电子电导 & 来自左侧（电正性）一侧的金属——Na、K、Ca——配上右侧元素——O、F、Cl，最好还是低原子量的元素 \\
金属 & Na、Al & 金属键——电子因在金属中可自由运动而动能低 & 结合力中等至强，密堆积；$R$ 低，$\mathrm{d}R/\mathrm{d}T>0$（金属导电）。金属光泽（即低频吸收带） & 出现最广：位于 H、B、Si、As、Te 以左者；若其右侧没有原子，它是金属；但即便有，往往也仍是金属。元素的随机组合通常是金属 \\
价晶体 & 金刚石 & “价键”——类似有机化学中的键。在原理上有别于金属？ & 硬，结合强，疏松堆积结构。$n^2$ 与 $\epsilon$ 高。$R$ 可变，但导电迁移率高 & 少见但重要。低 $Z$ 的居中元素：C、N、P、Si、B 为典型组分 \\
\bottomrule
\end{tabular}
\end{table}

在每一类晶体名下，我列出了其最显著的性质；这些性质在几乎所有情形下都可以证明是由关于晶体结合力的观念得出的。我希望在这些讲义的进程中，这些联系能在你们面前变得清楚起来。最后，把每一类别在周期表中出现的区域作为最后一列写下来，也是有意义的。

关于这张出现区域表，有一件值得注意却不常被讨论的事。如果我们只看周期表左右两侧的元素：像 Li、Na 这样在 s 与 p 带中多出一个电子的元素，乃至只有一个 p 电子的 Al，与那些 p 带中有一、两个甚至三个空穴的元素相对照，就会发现前者是金属，而后者倾向于形成分子晶体，或至多形成价晶体。人们常常被引导去相信空穴与电子之间存在一种基本的对称性，因此值得指出：在化学上它们颇为不同。往后在讨论磁性时，我们将在 d 壳层中少数空穴的行为与少数电子的行为之间发现颇为同样的二分性。这两类现象可能有联系，我们将在后面讨论。

同样重要的是讨论这种分类在多大程度上会失效。自然，人们立刻可以想到许多分类显然失效的例子，但其原则仍然成立——例如卤化铵 $(\mathrm{NH}_4)^+\mathrm{Cl}^-$ 等：这里 $\mathrm{NH}_4^+$ 是分子性结合的，而晶体整体是离子性的，并带有轻微的氢键色彩。分子式样的基团出现在离子晶体中是相当常见的。

远为重要也远为有趣的，是存在一大批物质——如今人们已经明白，它们把金属、离子晶体与价晶体这三个强结合类别之间的疆域完全填满了。例如，从 NaCl 出发——它在本性上几乎纯粹是离子型的——在周期表同一行中逐步提高两种组分的化合价。MgS 仍然停留在典型的离子型 NaCl 结构中，并遵守人们对离子晶体预期的绝大部分关系。AlP 则不然，它已经属于具有 ZnS 结构（闪锌矿型，zincblende）的 III--V 族半导体性类价晶体之列；研究得更为透彻的是与之极为相近的 InP 与 AlSb。我们都知道，Si 本身是卓越的价半导体（valence semiconductor）。随着原子序数增大，真正的离子晶体变得越来越罕见；ZnS 本身、CdS、CdTe 都表现出很强的价晶体乃至金属特征。另一方面，Ca、Sr、Ba 则倾向于保持其形成类离子氧化物与硫化物的性质——尽管人们本可以借助价晶体的论证，例如把这一点解释为：对它们而言下一个可用的轨道（即在其近邻 Sc 与 V 中被填充的那些轨道）是 d 轨道而不是 p 轨道，因而不适合四面体成键。已经有人定量地证明，大多数碱金属卤化物是相当好的离子晶体，其意义是两种离子上的电荷相当接近 $\pm 1$；但在其他方面，即便对氧化物而言，情况究竟如此到何种定量程度，也大可怀疑。例如硅酸盐，就典型地结晶成类价的而非离子型的结构。

更令人不安的是这样一个事实：随着我们对金属间化合物（intermetallic compound）这一广阔而少有研究的领域认识的加深，价晶体与金属之间、离子晶体与金属之间的区分正逐渐失去其泾渭分明的性质。我不指望能让你们对这一领域获得切实的洞察，因为对于缤纷现象背后的量子化学，几乎没有什么清晰的观念。一个例子即可说明事情能有多么糟糕，那就是化学计量的金属间化合物 NaTl (8)。这是一种导电良好的金属，不是半金属（semimetal）；其结构是 Tl 原子构成一个金刚石晶格，Na 原子占据 Tl 晶格中的大间隙位置，而这些间隙位置本身又构成一个金刚石晶格。既然 Tl 是如此之大的一个离子，这种结构绝非密堆积；事实上，它正是典型的价键型晶体才会具有的那种结构。要理解这件事，唯一的办法是假定 Na 已电离为 $\mathrm{Na}^+$，把那个多余的电子交给 Tl，使之成为 $\mathrm{Tl}^-$，从而具有适合金刚石晶格的 sp$^3$ 组态。Na 确实被电离了这一点，已由核共振证据所证实，其细节我在这里不打算讨论。于是，在单一物质中我们同时看到了数量级合理的金属导电、离子性的电荷转移，以及一种开放的、价键型的结构。

以上就是对固体量子化学一般特征的简要概述。我愿意重申：我并未追求完备，即便相对于 Seitz 的第一章也是如此；比如他在那里讨论了有序—无序、原子尺寸关系与 Hume-Rothery 合金，我略去这些论题只是因为时间不够，也因为文献中已有很好的处理，而不是因为它们不重要或乏味。我要强调的是，整整一个领域——金属间化合物——几乎仍未被开发。

界线——至少在价晶体与金属之间的界线——开始变得模糊的另一种途径，是人们普遍认识到：像 Si 这样可以是一种相当好的绝缘体的价晶体，其基本电子结构若从广义上看，与良金属的电子结构并无显著差别。就是说，当人们用诸如等离体共振（plasma resonance）或正电子湮没（positron annihilation）这类工具来研究电子时——这类工具只能在伏特量级、即十分之几个 Rydberg 的尺度上，而不是在十分之一伏特的尺度上分辨电子密度分布等等——这些电子看起来非常接近自由电子的费米球（Fermi sphere）；对许多良金属来说也是如此。必须强调，这类晶体中结合力的本性迄今仍无法付诸计算，因此对于结合能究竟来自何处、它们究竟更是价键型的还是金属型的，我们并没有一个清晰的定量概念。

还可以在这里对固体的另外几种分类方式作一番粗略的讨论。例如，有它们的磁性质的问题。绝大多数固体要么是抗磁性的，要么——对相当多的金属而言——是温度依赖性很弱的轻微顺磁性。然而，有几类物质，有证据表明其中的原子在固态仍保有自由的、可取向的自旋，有时还有轨道磁矩。这可以由强烈依赖温度且在低温下增大的顺磁性的存在来证明，也可以由铁磁或反铁磁现象——即磁矩的有序排列，它们大多出现在相当低的温度下——来证明。

这类磁矩出现得最广的，也许是在具有未满 3d、4f 或 5f（其次 4d 与 5d）电子壳层的原子的多少带离子性的盐类中——即铁组、稀土组与锕系组。典型的例子是 $\mathrm{MnF}_2$。可以相当稳妥地概括说：这些是处于内壳层中的电子，其意义是它们的密度的主体位于外层价电子密度的最后一个极大值之内。离子晶体中这类磁矩的出现与否，遵循相当著名且已被很好理解的离子络合物量子化学规则。

另一个类别——磁性分子晶体，即有机自由基的固态晶体——虽然对化学家颇有兴趣，在这里却大概只值一个脚注。

另一方面，在不是那么离子性的晶体中——如 3d 组较早成员 Ti、V 的氧化物与硫化物，或 4d 元素的氧化物与硫化物——情况就远非如此分明了，Morin 在一系列研究 (9) 中已表明这一点；磁性状态与非磁性状态之间的边界仍在划定之中，尚未被完全理解。

金属中的情形就更加混乱。这里，除了 f 组之外，已知情形中只有两个例外，凡有磁性的物质都含有铁组较后成员——从 Mn 到 Ni——之一的元素。磁性元素构成的化合物中，有许多并不具有磁性（例如 $\mathrm{CoSi}_2$——它几乎是被发现的第一个由非超导元素组成的超导化合物）；而到目前为止，只发现两种不含磁性组分的金属 $\mathrm{ZrZn}_2$ 与 $\mathrm{ScIn}_3$ 具有磁性。几乎不必说，这种局面的量子化学在撰写本书之际几乎仍是一个完全的谜。

另一种现象——超导电性（superconductivity）——似乎给了我们一种关于金属元素、合金与金属间化合物\footnote{以下我将把“合金”（alloy）定义为组分比例可在某个可感知范围内变动、因而其原子排列想必至少是部分无规的物质；而把“化合物”定义为具有完全确定的化学计量比与原子排列的物质。为方便起见，对后者不使用“合金”一词。}的粗略分类，它相对而言很少依赖于结构细节，而在某种宽泛的意义上依赖于基础量子化学。在金属之中，超导电性看来是常规而非例外。事实上，人们可以这样陈述它出现的要求：在非磁性金属中，实际上唯一的不超导“孤岛”出现在价电子数对原子数之比非常小之处——通常为 2 或更小——以及第六列的少数几个元素附近——W（Mo 已不再属于此类）——还有价电子对原子数之比与之相近的合金。磁性金属与半金属也不超导。同样，这些规则的原因充其量只有定性层面的理解 (10)。
